\chapter{RDMA / RoCE 网络基础}
\label{chap:rdma-roce}

上一章讨论了单机多卡通信：在一台服务器内部，GPU 可以通过 PCIe、NVLink 与 NVSwitch 交换数据。但是，当模型规模、数据规模或并发请求规模继续扩大时，单机已经不够。训练一个千亿参数模型可能需要数百到数千张 GPU；服务一个高并发大模型系统可能需要跨越多个节点部署 prefill、decode、embedding、rerank、tool calling 或多模型实例。此时，通信路径从“GPU 到 GPU”进一步扩展为：

\[
  \text{GPU} \rightarrow \text{PCIe/NVLink} \rightarrow \text{NIC} \rightarrow \text{Switch Fabric} \rightarrow \text{NIC} \rightarrow \text{GPU}.
\]

这就是本章要讨论的问题：多机 GPU 集群中的网络如何影响训练和推理性能。

在传统 Web 服务中，网络通常承载 RPC、HTTP、数据库访问、日志、对象存储等流量。对这些系统而言，毫秒级延迟往往已经足够，吞吐和可用性更重要。而在大模型训练中，GPU 每一步训练都可能需要跨节点同步数 GB 到数十 GB 的梯度、参数分片或激活张量。GPU 计算速度越快，网络越容易成为瓶颈。如果网络跟不上，GPU 就会在通信 barrier 前等待，昂贵的 Tensor Core 将处于空转状态。

因此，AI 集群网络不是普通数据中心网络的简单放大版。它必须同时满足：

\begin{itemize}
  \item 高带宽：支持大规模 all-reduce、reduce-scatter、all-gather、all-to-all；
  \item 低延迟：降低小消息、pipeline stage、decode 通信和控制路径等待；
  \item 低 CPU 开销：避免 CPU 参与数据搬运，减少内核协议栈开销；
  \item 拓扑可预测：让 collective 算法能够充分利用网络结构；
  \item 拥塞可控：避免少量热点流量拖慢整个训练作业；
  \item 可观测：能够定位链路、交换机、队列、NIC、NCCL 和作业层面的瓶颈。
\end{itemize}

RDMA、InfiniBand、RoCE 与 GPUDirect RDMA 正是为这类问题服务的关键技术。本章会从工程视角解释它们的基本概念，而不是陷入协议细节。读完本章后，你应该能够理解为什么 AI 集群常常选择 InfiniBand 或 RoCE 网络，为什么 RoCE 需要关注 PFC、ECN 和 lossless Ethernet，为什么 GPUDirect RDMA 能显著改善跨节点 GPU 通信，以及如何从 NCCL、网卡计数器和网络拓扑角度排查训练扩展性问题。

\section{多机训练为什么依赖高性能网络}
\label{sec:why-high-performance-network}

多机训练的核心矛盾是：计算被分散到多张 GPU 上，但模型更新仍然需要在逻辑上保持一致。无论使用 DP、TP、PP、ZeRO 还是 MoE，只要跨越节点，就会产生网络通信。随着 GPU 数量增加，单卡计算时间可能下降，但通信路径变长、通信规模变大、同步等待增加，最终扩展效率会受到网络限制。

\subsection{跨节点 All-Reduce}
\label{subsec:cross-node-allreduce}

数据并行中最典型的跨节点通信是梯度 all-reduce。假设有 $N$ 张 GPU，每张 GPU 持有一份完整模型副本，并处理不同 mini-batch。反向传播后，每张 GPU 得到本地梯度 $\nabla \theta_i$。为了让所有副本保持一致，需要计算平均梯度：

\[
  \nabla \theta = \frac{1}{N}\sum_{i=1}^{N}\nabla \theta_i.
\]

这通常通过 all-reduce 完成。每个 rank 输入自己的梯度 buffer，通信结束后每个 rank 都得到 reduce 后的完整结果。对于一个参数量为 $P$ 的模型，如果梯度使用 BF16 或 FP16，每个参数的梯度约为 2 字节，则一次梯度同步的逻辑数据规模大约为：

\[
  2P\ \text{Bytes}.
\]

对于十亿参数模型，这是 GB 级数据；对于百亿参数模型，梯度同步规模会更大。虽然实际系统会使用 bucket、overlap、reduce-scatter、ZeRO 等方法降低峰值显存和改善重叠，但通信总量仍然不容忽视。

在 ring all-reduce 中，每个 rank 的通信量大致为：

\[
  V_{\text{ring}} \approx 2 \cdot \frac{N-1}{N} \cdot S,
\]

其中 $S$ 是待 all-reduce tensor 的大小。随着 $N$ 增大，单 rank 通信量接近 $2S$。如果网络有效带宽为 $B$，仅从带宽角度估计通信时间下界为：

\[
  T_{\text{comm}} \gtrsim \frac{2S}{B}.
\]

这个估计忽略了延迟、协议开销、拓扑、拥塞和实现细节，但足以说明一个问题：当梯度 buffer 很大时，网络带宽直接决定同步时间。

\begin{figure}[H]
  \centering
  \begin{tikzpicture}[
    font=\small,
    nodebox/.style={draw, rounded corners, minimum width=2.0cm, minimum height=1.0cm, align=center},
    gpu/.style={draw, rounded corners, minimum width=1.0cm, minimum height=0.55cm, align=center},
    sw/.style={draw, rounded corners, minimum width=2.0cm, minimum height=0.75cm, align=center, fill=gray!10},
    arrow/.style={-{Stealth[length=2mm]}, thick}
  ]
    \node[nodebox] (n0) at (-4,1.4) {Node 0};
    \node[gpu] (g00) at (-4.55,0.75) {G0};
    \node[gpu] (g01) at (-3.45,0.75) {G1};

    \node[nodebox] (n1) at (0,1.4) {Node 1};
    \node[gpu] (g10) at (-0.55,0.75) {G2};
    \node[gpu] (g11) at (0.55,0.75) {G3};

    \node[nodebox] (n2) at (4,1.4) {Node 2};
    \node[gpu] (g20) at (3.45,0.75) {G4};
    \node[gpu] (g21) at (4.55,0.75) {G5};

    \node[sw] (leaf0) at (-2,3.1) {Leaf};
    \node[sw] (leaf1) at (2,3.1) {Leaf};
    \node[sw] (spine) at (0,4.5) {Spine};

    \draw[arrow] (n0.north) -- (leaf0.south);
    \draw[arrow] (n1.north) -- (leaf0.south);
    \draw[arrow] (n1.north) -- (leaf1.south);
    \draw[arrow] (n2.north) -- (leaf1.south);
    \draw[arrow] (leaf0.north) -- (spine.south);
    \draw[arrow] (leaf1.north) -- (spine.south);

    \draw[arrow, dashed] (g00.south) .. controls (-2.2,-0.7) and (2.2,-0.7) .. (g21.south);
    \node[align=center] at (0,-1.15) {跨节点 collective：GPU 数据经 NIC 与交换网络在节点间移动};
  \end{tikzpicture}
  \caption{多节点训练通信路径：节点内 GPU 通过本地互联到达 NIC，节点间数据经过交换网络传输。}
  \label{fig:multi-node-training-path}
\end{figure}

\subsection{网络带宽与训练扩展性}
\label{subsec:network-bandwidth-scaling}

训练扩展性的目标是：增加 GPU 数量后，单位时间完成的训练 token 或 sample 尽可能线性增长。理想情况下，$N$ 张 GPU 的吞吐是单卡的 $N$ 倍。但现实中，扩展效率总是低于 100\%。一个简化的 step time 模型可以写成：

\[
  T_{\text{step}} = T_{\text{compute}} + T_{\text{comm}} - T_{\text{overlap}} + T_{\text{bubble}} + T_{\text{misc}}.
\]

其中：

\begin{itemize}
  \item $T_{\text{compute}}$ 是 GPU 计算时间；
  \item $T_{\text{comm}}$ 是通信时间；
  \item $T_{\text{overlap}}$ 是计算与通信重叠抵消掉的部分；
  \item $T_{\text{bubble}}$ 是 pipeline bubble、同步等待等空洞；
  \item $T_{\text{misc}}$ 包括数据加载、checkpoint、CPU 调度等其他开销。
\end{itemize}

网络优化的目标不是让 $T_{\text{comm}}$ 变成 0，而是让它尽量小，并尽可能与计算重叠。当 GPU 更快、模型更大、并行维度更多时，通信在 step time 中的占比会上升。尤其在下面几类场景中，网络更容易成为瓶颈：

\begin{itemize}
  \item 使用较小 micro-batch，单步计算时间短，通信难以隐藏；
  \item TP 或 EP 跨节点，通信频率高；
  \item ZeRO-3 每层 all-gather 参数，通信分散且频繁；
  \item MoE all-to-all 产生复杂 east-west traffic；
  \item 网络存在 oversubscription，多个作业互相争用带宽；
  \item RoCE 拥塞控制不当，出现 PFC storm、丢包、重传或长尾延迟。
\end{itemize}

\subsection{Latency-Sensitive 与 Bandwidth-Sensitive}
\label{subsec:latency-vs-bandwidth-sensitive}

不同通信模式对网络的要求不同。可以粗略分为 latency-sensitive 和 bandwidth-sensitive 两类。

\textbf{Latency-sensitive} 通信的特点是消息较小、频率较高、在关键路径上等待。例如：

\begin{itemize}
  \item pipeline parallel 相邻 stage 之间的小 activation 传输；
  \item 推理 decode 阶段每层 tensor parallel collective；
  \item 控制面元数据同步；
  \item 小 batch 或小 tensor 上的 all-reduce；
  \item MoE token dispatch 中大量小粒度消息。
\end{itemize}

这类通信即使数据量不大，也会因为等待和同步影响端到端延迟。优化重点是降低 RTT、减少协议栈开销、避免排队、减少 CPU 参与和减少不必要同步。

\textbf{Bandwidth-sensitive} 通信的特点是消息大、传输时间主要由链路带宽决定。例如：

\begin{itemize}
  \item 数据并行大 bucket 梯度 all-reduce；
  \item ZeRO reduce-scatter / all-gather 大参数分片；
  \item checkpoint 分布式保存；
  \item 大 batch embedding 或 parameter server 风格通信；
  \item 大规模 MoE all-to-all。
\end{itemize}

这类通信的优化重点是提高有效带宽、增加多 rail 并行、选择合适 collective 算法、避免 oversubscription、让通信与计算重叠。

\begin{table}[H]
  \centering
  \caption{Latency-sensitive 与 Bandwidth-sensitive 通信对比}
  \label{tab:latency-bandwidth-communication}
  \begin{tabular}{p{3.2cm}p{4.6cm}p{4.6cm}p{3.0cm}}
    \toprule
    类型 & 特征 & 典型场景 & 优化重点 \\
    \midrule
    Latency-sensitive & 小消息、高频、同步等待明显 & Decode TP、PP stage、控制面同步、小 tensor collective & 降低 RTT、减少 CPU 参与、减少同步和排队 \\
    Bandwidth-sensitive & 大消息、传输时间由带宽主导 & DP all-reduce、ZeRO、checkpoint、MoE 大规模 all-to-all & 提高链路利用率、多 rail、拓扑感知 collective \\
    Mixed & 消息大小变化大，既受延迟也受带宽影响 & MoE token dispatch、动态 batch serving、参数分片加载 & 分层调度、分桶、合并小消息、拥塞控制 \\
    \bottomrule
  \end{tabular}
\end{table}

\subsection{通信瓶颈如何吞噬 GPU 利用率}
\label{subsec:communication-bottleneck-gpu-utilization}

通信瓶颈经常表现为 GPU 利用率下降，但它的根源并不在 GPU kernel 本身。一个训练 step 中，GPU 可能先执行 forward 和 backward kernel，然后等待 NCCL collective 完成，再进入下一层或下一步。如果通信无法与计算重叠，timeline 上就会出现明显空洞。

更隐蔽的情况是：NCCL kernel 正在 GPU 上运行，但 SM utilization 并不高。此时从 \texttt{nvidia-smi} 看 GPU 可能不是完全空闲，但 Tensor Core 没有执行模型计算。训练吞吐仍然下降。对于 AI Infra 工程师，关键不是追求某个单一利用率指标，而是判断 GPU 时间到底花在：

\begin{itemize}
  \item 计算 kernel；
  \item NCCL 通信 kernel；
  \item memory copy；
  \item CPU launch 等待；
  \item stream synchronization；
  \item dataloader 或存储等待；
  \item checkpoint I/O。
\end{itemize}

这需要结合 Nsight Systems、NCCL 日志、网络计数器和训练框架日志一起看。一个成熟的性能分析流程一定是跨层的：从模型 step time 到 NCCL collective，从 GPU timeline 到 NIC counters，从 rank mapping 到交换机队列。

\section{RDMA 基础}
\label{sec:rdma-basics}

RDMA 是 Remote Direct Memory Access 的缩写，即远程直接内存访问。它允许一台机器上的网卡直接读写另一台机器上的内存区域，而不需要远端 CPU 频繁参与数据拷贝，也不需要每个数据包都经过传统内核网络协议栈的完整处理路径。

RDMA 的核心价值可以概括为三点：

\begin{enumerate}
  \item \textbf{低延迟}：绕过传统内核协议栈，减少系统调用、协议处理和上下文切换；
  \item \textbf{低 CPU 开销}：数据搬运由 NIC 和 DMA 引擎完成，CPU 不需要参与每次拷贝；
  \item \textbf{高带宽}：更容易接近网卡和链路的有效带宽上限。
\end{enumerate}

在 AI 集群中，RDMA 通常不是直接由训练代码手写 verbs API，而是被 NCCL、UCX、MPI、SHARP、NVSHMEM 或框架通信层间接使用。尽管如此，理解 RDMA 的基本对象仍然非常有帮助，因为很多故障和性能问题会暴露为 QP、CQ、MR、GID、LID、MTU、PFC、ECN、RoCE mode 等配置问题。

\subsection{Kernel Bypass}
\label{subsec:kernel-bypass}

传统 TCP/IP 数据路径通常需要经过操作系统内核协议栈。发送数据时，应用把数据交给内核，内核处理 socket、TCP/IP 协议、拥塞控制、拷贝、网卡队列等；接收数据时，数据包到达网卡后，内核处理中断、协议栈、缓冲区，再把数据交给应用。这个路径通用、可靠、易用，但对高频低延迟通信来说开销较大。

RDMA 的 kernel bypass 思想是：连接建立、资源注册等控制操作仍然需要内核参与，但数据路径尽量绕过内核。应用通过用户态库把 work request 提交给网卡队列，网卡直接从已注册内存读取或写入数据，并在完成后通过 completion queue 通知应用。

\begin{figure}[H]
  \centering
  \begin{tikzpicture}[
    font=\small,
    box/.style={draw, rounded corners, minimum width=2.3cm, minimum height=0.7cm, align=center},
    wide/.style={draw, rounded corners, minimum width=3.2cm, minimum height=0.75cm, align=center},
    arrow/.style={-{Stealth[length=2mm]}, thick}
  ]
    \node[wide] (app1) at (-4,4.0) {Application};
    \node[wide] (kernel1) at (-4,2.8) {Kernel TCP/IP};
    \node[wide] (nic1) at (-4,1.6) {NIC};
    \node[wide] (net1) at (-4,0.4) {Network};

    \draw[arrow] (app1) -- (kernel1);
    \draw[arrow] (kernel1) -- (nic1);
    \draw[arrow] (nic1) -- (net1);

    \node[wide] (app2) at (3.2,4.0) {Application};
    \node[wide] (verbs) at (3.2,2.8) {User-space RDMA Lib};
    \node[wide] (nic2) at (3.2,1.6) {RNIC / HCA};
    \node[wide] (net2) at (3.2,0.4) {Network};

    \draw[arrow] (app2) -- (verbs);
    \draw[arrow] (verbs) -- (nic2);
    \draw[arrow] (nic2) -- (net2);

    \node[box, minimum width=2.4cm] at (-4,-0.8) {传统 TCP 路径};
    \node[box, minimum width=2.4cm] at (3.2,-0.8) {RDMA 数据路径};

    \draw[dashed] (-0.6,-1.3) -- (-0.6,4.5);
    \node[align=center] at (-0.6,4.85) {对比};
  \end{tikzpicture}
  \caption{传统 TCP 与 RDMA 数据路径对比：RDMA 的数据面尽量绕过内核协议栈。}
  \label{fig:tcp-vs-rdma-path}
\end{figure}

Kernel bypass 不是说内核完全不参与。RDMA 仍然需要内核驱动管理设备、权限、内存注册、连接建立和安全隔离。所谓 bypass 主要指高频数据收发路径不再为每个消息经过传统 TCP/IP 内核处理。

\subsection{Zero-Copy}
\label{subsec:zero-copy}

Zero-copy 指数据在发送和接收过程中尽量避免 CPU 参与的额外内存拷贝。传统网络通信中，数据可能经历应用缓冲区到内核缓冲区、内核缓冲区到网卡 DMA buffer 等多次移动。RDMA 中，应用先把一段内存注册给 RDMA 设备，NIC 获得访问这段内存的权限和地址映射。之后，NIC 可以直接通过 DMA 从这段内存读取或写入。

在 AI 训练中，zero-copy 的意义非常大。梯度、参数分片、activation 或 KV Cache 都可能是大 tensor。如果每次跨节点通信都需要 CPU 拷贝一份，CPU 内存带宽、缓存污染和延迟都会变成负担。RDMA 让通信库更接近“从源 buffer 直接送到目标 buffer”的理想路径。

不过 zero-copy 也不是免费的。为了让 NIC 安全、正确地访问内存，需要 memory registration。注册过程可能涉及 pin memory、建立地址转换、设置访问权限等操作，成本不低。因此高性能通信库通常会复用注册内存，避免频繁注册和注销。

\subsection{Queue Pair}
\label{subsec:queue-pair}

Queue Pair，简称 QP，是 RDMA 编程模型中的核心对象。一个 QP 通常包含：

\begin{itemize}
  \item Send Queue：应用向其中提交发送、写、读等 work request；
  \item Receive Queue：用于接收远端发送过来的消息；
  \item 与之关联的 Completion Queue：用于获取完成事件。
\end{itemize}

可以把 QP 理解为 RDMA 连接或通信端点的抽象。不同 QP 类型支持不同通信语义，例如可靠连接、非可靠连接、数据报等。AI 集群中最常见的上层用户通常不直接管理 QP，但 NCCL、UCX、MPI 等通信库会在底层创建并管理这些对象。

\begin{figure}[H]
  \centering
  \begin{tikzpicture}[
    font=\small,
    box/.style={draw, rounded corners, minimum width=2.6cm, minimum height=0.7cm, align=center},
    host/.style={draw, rounded corners, minimum width=4.0cm, minimum height=4.4cm, align=center},
    arrow/.style={-{Stealth[length=2mm]}, thick}
  ]
    \node[host] (h0) at (-3.2,1.8) {};
    \node[anchor=north west] at ($(h0.north west)+(0.2,-0.2)$) {Host A};
    \node[box] (sq0) at (-3.2,2.8) {Send Queue};
    \node[box] (rq0) at (-3.2,1.8) {Receive Queue};
    \node[box] (cq0) at (-3.2,0.8) {Completion Queue};

    \node[host] (h1) at (3.2,1.8) {};
    \node[anchor=north west] at ($(h1.north west)+(0.2,-0.2)$) {Host B};
    \node[box] (sq1) at (3.2,2.8) {Send Queue};
    \node[box] (rq1) at (3.2,1.8) {Receive Queue};
    \node[box] (cq1) at (3.2,0.8) {Completion Queue};

    \draw[arrow] (sq0.east) -- node[above] {Work Request} (rq1.west);
    \draw[arrow] (sq1.west) -- node[below] {Work Request} (rq0.east);
    \draw[arrow] (rq0.south) -- (cq0.north);
    \draw[arrow] (rq1.south) -- (cq1.north);

    \node[align=center] at (0,-0.6) {QP 将用户态提交的 work request 与网卡数据路径连接起来};
  \end{tikzpicture}
  \caption{RDMA Queue Pair 抽象：发送队列、接收队列和完成队列共同描述通信端点。}
  \label{fig:rdma-queue-pair}
\end{figure}

\subsection{Completion Queue}
\label{subsec:completion-queue}

Completion Queue，简称 CQ，用来通知应用某个 work request 已经完成。应用可以轮询 CQ，也可以使用事件通知。高性能通信库通常偏向轮询，因为轮询可以避免中断和上下文切换开销，但会消耗 CPU core。许多 AI 集群会为通信线程、NCCL proxy、MPI progress engine 等预留 CPU 资源，以免这些通信路径被业务线程或系统任务干扰。

CQ 的存在提醒我们：RDMA 降低了 CPU 数据搬运开销，但并不意味着 CPU 完全无关。连接管理、队列推进、completion 处理、错误处理、通信库调度仍然需要 CPU 或专门的硬件机制支持。如果 CPU 过载，网络性能仍然可能下降。

\subsection{Memory Registration}
\label{subsec:memory-registration}

RDMA 设备不能随意访问进程的任意虚拟内存。应用需要先注册内存区域，得到 Memory Region，简称 MR。注册后，设备获得这段内存的访问权限和地址映射，并生成 local key 或 remote key，用于本地或远端访问验证。

Memory registration 的工程意义包括：

\begin{itemize}
  \item 注册内存通常需要 pin 住物理页，防止被操作系统换出或移动；
  \item 注册和注销有成本，不应在热路径中频繁执行；
  \item 通信库通常会维护 registration cache；
  \item 容器、权限、IOMMU、驱动和内核模块都可能影响注册行为；
  \item GPUDirect RDMA 中，注册对象可能是 GPU memory，而不只是 CPU memory。
\end{itemize}

对于大模型训练，频繁分配和释放通信 buffer 可能导致 registration cache miss、内存碎片和性能抖动。因此框架通常会复用 gradient bucket、通信 workspace、参数 shard buffer 等内存区域。

\section{InfiniBand 与 RoCE}
\label{sec:infiniband-and-roce}

RDMA 是一种能力或编程模型，而 InfiniBand 和 RoCE 是实现 RDMA 网络的两类常见方式。它们在 AI 集群中都很常见。

\subsection{InfiniBand 网络}
\label{subsec:infiniband-network}

InfiniBand 是为高性能计算和低延迟通信设计的专用网络技术。它通常由 HCA、InfiniBand switch、subnet manager、专用线缆和协议栈组成。InfiniBand 的优势在于端到端设计较为一致，天然支持 RDMA，延迟低、带宽高、拥塞控制和拓扑管理能力成熟。

在 AI 训练集群中，InfiniBand 常用于跨节点 GPU 通信。NCCL 可以通过 InfiniBand HCA 在节点间传输 collective 数据，并结合 GPUDirect RDMA 实现 GPU memory 与 NIC 的直接数据路径。

InfiniBand 的典型运维关注点包括：

\begin{itemize}
  \item HCA 速率、端口状态和 link width；
  \item subnet manager 是否稳定；
  \item LID、SL、MTU、GID 等配置；
  \item switch 端口错误计数；
  \item fabric 拓扑和路由；
  \item 多 rail 绑定和 rank 到 NIC 的映射；
  \item NCCL 是否正确选择 IB 设备。
\end{itemize}

\subsection{RoCEv1 / RoCEv2}
\label{subsec:rocev1-rocev2}

RoCE 是 RDMA over Converged Ethernet，即在以太网上承载 RDMA。它的吸引力在于可以复用以太网生态：交换机、布线、运维体系、IP 网络能力等。RoCE 主要有两个版本：

\begin{itemize}
  \item \textbf{RoCEv1}：运行在二层以太网上，不能跨三层路由；
  \item \textbf{RoCEv2}：把 RDMA transport 封装在 UDP/IP 上，可以跨三层网络。
\end{itemize}

RoCEv2 是现代数据中心更常见的形态，因为它更容易与 leaf-spine、L3 routing、ECMP 等数据中心网络设计结合。但 RoCE 的挑战也很明显：传统以太网默认不是无损网络，而 RDMA 对丢包非常敏感。如果发生丢包，性能可能急剧下降。因此 RoCE 网络通常需要配置 PFC、ECN、QoS、buffer、拥塞控制和优先级映射。

\begin{figure}[H]
  \centering
  \begin{tikzpicture}[
    font=\small,
    layer/.style={draw, rounded corners, minimum width=8.8cm, minimum height=0.65cm, align=center},
    arrow/.style={-{Stealth[length=2mm]}, thick}
  ]
    \node[layer] (app) at (0,4.2) {NCCL / UCX / MPI / Application};
    \node[layer] (verbs) at (0,3.35) {RDMA Verbs / Communication Runtime};
    \node[layer] (ibtransport) at (0,2.5) {InfiniBand Transport};
    \node[layer] (udp) at (0,1.65) {UDP / IP（RoCEv2）};
    \node[layer] (eth) at (0,0.8) {Ethernet + QoS / PFC / ECN};
    \node[layer] (phy) at (0,-0.05) {NIC / Switch / Physical Link};

    \draw[arrow] (app) -- (verbs);
    \draw[arrow] (verbs) -- (ibtransport);
    \draw[arrow] (ibtransport) -- (udp);
    \draw[arrow] (udp) -- (eth);
    \draw[arrow] (eth) -- (phy);

    \node[align=center] at (5.4,2.1) {RoCEv2：\\将 RDMA 传输封装\\在 UDP/IP/Ethernet 上};
  \end{tikzpicture}
  \caption{RoCEv2 协议栈抽象：RDMA transport 通过 UDP/IP 承载在以太网 fabric 上。}
  \label{fig:roce-v2-stack}
\end{figure}

\subsection{Lossless Ethernet}
\label{subsec:lossless-ethernet}

RoCE 的关键工程目标是构造近似无损的以太网环境。所谓 lossless 并不是说永远不会出错，而是通过链路层流控、拥塞通知、队列管理和 QoS 配置，尽量避免因拥塞丢包导致 RDMA 重传或连接异常。

在普通 TCP 网络中，丢包常被用作拥塞信号，TCP 会重传并降低发送速率。但 RDMA transport 对丢包更敏感，重传和恢复成本高。对于大规模 all-reduce 或 all-to-all，少量链路丢包可能造成整个训练 step 长尾变长，因为 collective 通常需要等待最慢 rank。

Lossless Ethernet 的核心组件包括：

\begin{itemize}
  \item \textbf{PFC}：Priority Flow Control，在特定优先级上暂停上游发送，避免队列溢出；
  \item \textbf{ECN}：Explicit Congestion Notification，在不丢包的情况下标记拥塞；
  \item \textbf{DCQCN 或类似拥塞控制}：发送端根据拥塞信号调节速率；
  \item \textbf{QoS / DSCP / PCP 映射}：确保 RoCE 流量进入正确优先级队列；
  \item \textbf{交换机 buffer 配置}：避免瞬时突发导致丢包；
  \item \textbf{ECMP 与路由设计}：在多路径网络中平衡流量。
\end{itemize}

\subsection{PFC 与 ECN}
\label{subsec:pfc-ecn}

PFC 和 ECN 经常一起出现，但它们解决的问题不同。

PFC 是链路层暂停机制。当某个交换机端口上的特定优先级队列接近拥塞阈值时，可以向上游发送 pause frame，要求对方暂停该优先级流量。这样可以避免队列溢出丢包。它的优点是直接有效，缺点是可能造成 head-of-line blocking、拥塞扩散，甚至在配置不当时出现 PFC storm。

ECN 是网络层拥塞标记机制。当交换机检测到队列拥塞时，不直接丢包，而是在 IP header 中标记 ECN。接收端或协议栈把拥塞信号反馈给发送端，发送端降低速率。相比 PFC，ECN 更适合做端到端拥塞控制，但它要求网络路径上的设备都正确支持和配置。

\begin{table}[H]
  \centering
  \caption{PFC 与 ECN 的工程对比}
  \label{tab:pfc-ecn-comparison}
  \begin{tabular}{p{2.8cm}p{5.0cm}p{5.0cm}}
    \toprule
    机制 & 优势 & 风险 \\
    \midrule
    PFC & 可以快速阻止特定优先级队列丢包，适合构造无损队列 & 可能导致 head-of-line blocking、拥塞扩散、pause storm \\
    ECN & 不丢包即可反馈拥塞，支持端到端速率调节 & 需要端到端设备和主机协议栈正确配置，阈值不当会影响吞吐或延迟 \\
    QoS 映射 & 将 RoCE 流量放入专用优先级，避免与普通流量互相干扰 & 映射错误会让 RoCE 流量进入非无损队列或错误拥塞控制域 \\
    Buffer 调优 & 吸收突发流量，降低瞬时丢包概率 & 过大 buffer 可能增加排队延迟，过小 buffer 容易丢包或触发 PFC \\
    \bottomrule
  \end{tabular}
\end{table}

AI 集群中的 PFC/ECN 配置不能照搬普通业务网络。训练流量具有明显 collective 特征：all-reduce、reduce-scatter、all-gather、all-to-all 会产生同步、突发和多对多模式。错误的阈值可能在小规模测试中看不出问题，但在大规模训练中导致长尾延迟、NCCL timeout 或吞吐抖动。

\subsection{拥塞控制}
\label{subsec:congestion-control}

拥塞控制的目标是让网络在高负载下仍然稳定运行。对 RoCE 来说，拥塞控制尤其重要，因为它既要避免丢包，又要避免过度 PFC 导致拥塞扩散。

AI 集群中的拥塞来源包括：

\begin{itemize}
  \item 多个训练作业共享同一 leaf-spine 网络；
  \item MoE all-to-all 造成多对多突发；
  \item checkpoint 或数据加载与训练通信混用网络；
  \item 错误 rank mapping 导致某些链路热点；
  \item ECMP hash 不均匀，多个大流落到同一路径；
  \item oversubscription 过高；
  \item 某个节点或 NIC 异常导致重传和尾延迟。
\end{itemize}

拥塞控制不是只调一个参数。它涉及主机、NIC、交换机、QoS、队列、路由、NCCL 算法和作业调度。工程上应建立一套闭环：压测、监控、告警、定位、回归。不要只在训练作业报错后才查看交换机丢包计数；更好的做法是在集群上线前就用 NCCL tests、ib\_write\_bw、perftest、RoCE counters、交换机 telemetry 做基线测试。

\section{GPUDirect RDMA}
\label{sec:gpudirect-rdma}

RDMA 让 NIC 可以直接访问主机内存，而 GPUDirect RDMA 进一步让 NIC 可以直接访问 GPU memory。对于多节点 GPU 训练，这是非常关键的能力。没有 GPUDirect RDMA 时，跨节点通信可能需要经过 CPU host memory staging：

\[
  \text{GPU HBM} \rightarrow \text{CPU Memory} \rightarrow \text{NIC} \rightarrow \text{Network}.
\]

有 GPUDirect RDMA 时，理想路径可以变成：

\[
  \text{GPU HBM} \rightarrow \text{NIC} \rightarrow \text{Network}.
\]

接收方向同理，远端数据可以由 NIC 直接写入 GPU memory。这样可以减少 CPU 内存拷贝、降低延迟、释放 CPU 资源，并提高跨节点 GPU collective 的效率。

\subsection{GPU Memory 直接被 NIC 访问}
\label{subsec:nic-access-gpu-memory}

GPUDirect RDMA 需要 GPU、NIC、驱动、内核模块、PCIe 拓扑和系统配置共同支持。NIC 要能够通过 PCIe peer-to-peer DMA 访问 GPU memory；GPU 驱动需要暴露相应的 peer memory 支持；通信库需要正确注册 GPU buffer 并使用支持 GPUDirect 的数据路径。

从应用层看，这些细节通常被 NCCL 或 UCX 封装。例如 PyTorch 中一个跨节点 \texttt{dist.all\_reduce} 操作，底层可能通过 NCCL 使用 InfiniBand 或 RoCE，并利用 GPUDirect RDMA 直接传输 CUDA tensor。

\begin{figure}[H]
  \centering
  \begin{tikzpicture}[
    font=\small,
    box/.style={draw, rounded corners, minimum width=2.1cm, minimum height=0.75cm, align=center},
    nodebox/.style={draw, rounded corners, minimum width=5.0cm, minimum height=3.5cm, align=center},
    arrow/.style={-{Stealth[length=2mm]}, thick}
  ]
    \node[nodebox] (nodea) at (-3.3,1.8) {};
    \node[anchor=north west] at ($(nodea.north west)+(0.2,-0.2)$) {Node A};
    \node[box] (gpua) at (-4.2,2.4) {GPU HBM};
    \node[box] (nica) at (-2.4,1.2) {NIC / HCA};
    \node[box] (cpua) at (-4.2,1.2) {CPU Memory};

    \node[nodebox] (nodeb) at (3.3,1.8) {};
    \node[anchor=north west] at ($(nodeb.north west)+(0.2,-0.2)$) {Node B};
    \node[box] (gpub) at (4.2,2.4) {GPU HBM};
    \node[box] (nicb) at (2.4,1.2) {NIC / HCA};
    \node[box] (cpub) at (4.2,1.2) {CPU Memory};

    \node[box, minimum width=2.8cm] (net) at (0,1.2) {RDMA Network};

    \draw[arrow] (gpua) -- node[above left] {\scriptsize P2P DMA} (nica);
    \draw[arrow] (nica) -- (net);
    \draw[arrow] (net) -- (nicb);
    \draw[arrow] (nicb) -- node[above right] {\scriptsize P2P DMA} (gpub);

    \draw[dashed] (gpua) -- (cpua);
    \draw[dashed] (cpub) -- (gpub);

    \node[align=center] at (0,-0.3) {实线：GPUDirect RDMA 数据路径；虚线：传统 CPU staging 路径};
  \end{tikzpicture}
  \caption{GPUDirect RDMA 数据流：NIC 可以直接读写 GPU memory，减少 CPU host memory 中转。}
  \label{fig:gpudirect-rdma-dataflow}
\end{figure}

\subsection{绕过 CPU Copy}
\label{subsec:bypass-cpu-copy}

绕过 CPU copy 的价值不仅是减少一次数据拷贝。它还减少了多个隐性开销：

\begin{itemize}
  \item CPU 内存带宽不再成为跨节点 GPU 通信的中间瓶颈；
  \item CPU cache 不会被大 tensor 通信污染；
  \item CPU core 可以用于 dataloader、通信进度、调度和服务逻辑；
  \item 降低端到端延迟，尤其对小消息和频繁通信有帮助；
  \item 减少内存分配、pin memory 和 staging buffer 管理复杂度。
\end{itemize}

不过 GPUDirect RDMA 也有条件和限制。它依赖设备拓扑。如果 GPU 和 NIC 之间跨越不支持 P2P 的 PCIe switch 或平台配置不当，性能可能下降或功能不可用。多 NUMA、多 NIC、多 GPU 节点中，GPU 与 NIC 的亲和性尤其重要。上一章中 \texttt{nvidia-smi topo -m} 的 GPU-NIC 路径，在跨节点通信中会变得更加关键。

\subsection{NCCL 跨节点通信}
\label{subsec:nccl-cross-node-communication}

NCCL 在多节点环境中会自动选择可用的网络传输路径。常见环境变量包括：

\begin{itemize}
  \item \texttt{NCCL\_DEBUG}：打开调试日志；
  \item \texttt{NCCL\_DEBUG\_SUBSYS}：指定调试子系统；
  \item \texttt{NCCL\_SOCKET\_IFNAME}：指定 socket 网络接口；
  \item \texttt{NCCL\_IB\_HCA}：指定 InfiniBand / RoCE HCA；
  \item \texttt{NCCL\_IB\_GID\_INDEX}：选择 RoCE GID；
  \item \texttt{NCCL\_IB\_ROCE\_VERSION\_NUM}：指定 RoCE version；
  \item \texttt{NCCL\_NET\_GDR\_LEVEL}：影响 GPUDirect RDMA 使用策略。
\end{itemize}

下面是一个用于调试的多节点 NCCL 测试示例：

\begin{lstlisting}[language=bash,caption={多节点 NCCL 通信调试示例},label={lst:nccl-rdma-debug}]
export NCCL_DEBUG=INFO
export NCCL_DEBUG_SUBSYS=INIT,NET,GRAPH,TOPO,COLL
export NCCL_IB_HCA=mlx5_0,mlx5_1
export NCCL_SOCKET_IFNAME=eth0

mpirun -np 16 -H node0:8,node1:8 \
  ./build/all_reduce_perf -b 8M -e 8G -f 2 -g 1
\end{lstlisting}

NCCL 日志中需要重点观察：

\begin{itemize}
  \item 是否检测到 IB/RoCE 设备；
  \item 是否使用了预期 HCA；
  \item 是否启用 GPUDirect RDMA；
  \item channel 数量是否合理；
  \item ring/tree 拓扑是否符合节点结构；
  \item 是否 fallback 到 socket；
  \item 是否存在连接错误、timeout 或 retry。
\end{itemize}

如果 NCCL 意外 fallback 到 TCP socket，训练仍然可能运行，但性能会显著下降。很多“训练突然变慢”的问题，最终都能在 NCCL 初始化日志中发现线索。

\subsection{多 NIC、多 Rail 与拓扑感知}
\label{subsec:multi-nic-multirail}

高端 AI 服务器通常配备多张 NIC 或多端口 HCA。多 rail 的目标是让多个网络端口并行承载通信，提高总带宽并改善拓扑利用率。例如，一个 8 GPU 节点可能有 4 张 400Gb/s NIC，每两张 GPU 共享或靠近一张 NIC。理想情况下，rank 应该优先使用距离自己更近的 NIC，避免跨 NUMA 或跨 PCIe switch 访问远端 NIC。

多 rail 设计涉及三个映射：

\begin{enumerate}
  \item GPU 到 NIC 的物理拓扑映射；
  \item rank 到 GPU 的映射；
  \item NCCL channel 到 NIC rail 的映射。
\end{enumerate}

如果这些映射不一致，可能出现某些 NIC 过载、某些 NIC 空闲，或者 GPU 访问远端 NIC 导致 PCIe/NUMA 瓶颈。

\begin{figure}[H]
  \centering
  \begin{tikzpicture}[
    font=\small,
    gpu/.style={draw, rounded corners, minimum width=0.9cm, minimum height=0.55cm, align=center},
    nic/.style={draw, rounded corners, minimum width=1.2cm, minimum height=0.6cm, align=center, fill=gray!10},
    sw/.style={draw, rounded corners, minimum width=2.0cm, minimum height=0.7cm, align=center},
    fast/.style={line width=1.1pt},
    arrow/.style={-{Stealth[length=2mm]}, thick}
  ]
    \foreach \i/\x in {0/-3.5,1/-2.5,2/-1.5,3/-0.5,4/0.5,5/1.5,6/2.5,7/3.5} {
      \node[gpu] (g\i) at (\x,0.4) {G\i};
    }

    \node[nic] (n0) at (-3.0,2.0) {NIC0};
    \node[nic] (n1) at (-1.0,2.0) {NIC1};
    \node[nic] (n2) at (1.0,2.0) {NIC2};
    \node[nic] (n3) at (3.0,2.0) {NIC3};

    \node[sw] (fabric) at (0,3.6) {Network Fabric};

    \draw[fast] (g0.north) -- (n0.south);
    \draw[fast] (g1.north) -- (n0.south);
    \draw[fast] (g2.north) -- (n1.south);
    \draw[fast] (g3.north) -- (n1.south);
    \draw[fast] (g4.north) -- (n2.south);
    \draw[fast] (g5.north) -- (n2.south);
    \draw[fast] (g6.north) -- (n3.south);
    \draw[fast] (g7.north) -- (n3.south);

    \draw[arrow] (n0.north) -- (fabric.south);
    \draw[arrow] (n1.north) -- (fabric.south);
    \draw[arrow] (n2.north) -- (fabric.south);
    \draw[arrow] (n3.north) -- (fabric.south);

    \node[align=center] at (0,-0.75) {多 rail 节点：rank、GPU 与 NIC 的亲和性会影响跨节点通信性能};
  \end{tikzpicture}
  \caption{多 NIC / 多 rail 拓扑感知：GPU 应尽量使用拓扑距离更近的 NIC。}
  \label{fig:multi-nic-multirail}
\end{figure}

\section{AI 集群网络设计}
\label{sec:ai-cluster-network-design}

理解 RDMA 和 RoCE 之后，还需要从集群整体设计角度看网络。AI 集群网络不只是“每台机器插几张网卡”，还包括拓扑、交换机层级、oversubscription、路由、拥塞控制、监控和调度协同。

\subsection{Fat-Tree}
\label{subsec:fat-tree}

Fat-tree 是数据中心和 HPC 中常见的拓扑思想。它通过多层交换机构造高带宽 bisection bandwidth，让任意两个节点之间都有较好的通信路径。典型 leaf-spine 架构可以看作 fat-tree 的一种数据中心实现：服务器连接到 leaf switch，leaf switch 上联到 spine switch，spine 之间提供横向互联能力。

对于 AI 训练，fat-tree 的优点是拓扑规则、扩展性好、路径多，适合集体通信。缺点是成本高，尤其当目标是低 oversubscription 或 non-blocking fabric 时，需要大量高速交换机端口和光模块。

\begin{figure}[H]
  \centering
  \begin{tikzpicture}[
    font=\small,
    srv/.style={draw, rounded corners, minimum width=1.1cm, minimum height=0.55cm, align=center},
    sw/.style={draw, rounded corners, minimum width=1.5cm, minimum height=0.65cm, align=center, fill=gray!10},
    link/.style={line width=0.8pt}
  ]
    \node[sw] (s0) at (-3,4.0) {Spine0};
    \node[sw] (s1) at (-1,4.0) {Spine1};
    \node[sw] (s2) at (1,4.0) {Spine2};
    \node[sw] (s3) at (3,4.0) {Spine3};

    \node[sw] (l0) at (-3,2.2) {Leaf0};
    \node[sw] (l1) at (-1,2.2) {Leaf1};
    \node[sw] (l2) at (1,2.2) {Leaf2};
    \node[sw] (l3) at (3,2.2) {Leaf3};

    \foreach \l in {l0,l1,l2,l3} {
      \foreach \s in {s0,s1,s2,s3} {
        \draw[link] (\l.north) -- (\s.south);
      }
    }

    \foreach \i/\x in {0/-3.6,1/-2.4,2/-1.6,3/-0.4,4/0.4,5/1.6,6/2.4,7/3.6} {
      \node[srv] (n\i) at (\x,0.6) {Node\i};
    }

    \draw[link] (n0.north) -- (l0.south);
    \draw[link] (n1.north) -- (l0.south);
    \draw[link] (n2.north) -- (l1.south);
    \draw[link] (n3.north) -- (l1.south);
    \draw[link] (n4.north) -- (l2.south);
    \draw[link] (n5.north) -- (l2.south);
    \draw[link] (n6.north) -- (l3.south);
    \draw[link] (n7.north) -- (l3.south);

    \node[align=center] at (0,-0.35) {Leaf-Spine / Fat-tree：提供多路径、高横向带宽的 AI 集群网络};
  \end{tikzpicture}
  \caption{AI 集群 leaf-spine 网络拓扑示意：节点通过 leaf 接入，多 leaf 通过 spine 互联。}
  \label{fig:ai-cluster-leaf-spine}
\end{figure}

\subsection{Dragonfly}
\label{subsec:dragonfly}

Dragonfly 是另一类高性能网络拓扑，常见于大型 HPC 系统。它通过 group 内高带宽连接和 group 间全局连接，在较低直径下连接大量节点。相比 fat-tree，dragonfly 可以在某些规模下减少交换机和线缆成本，但路由、拥塞控制和流量工程更加复杂。

对于 AI 集群，dragonfly 的适用性取决于通信模式。如果 workload 以规则 collective 为主，并且通信库和路由能够充分利用 topology-aware 算法，dragonfly 可以提供良好性能。但如果多个作业混部、all-to-all 流量复杂、拥塞控制不足，热点路径可能导致长尾问题。

本书不深入展开 dragonfly 的路由算法，只强调一点：AI 集群网络拓扑必须与通信模式共同设计。大模型训练不是普通均匀流量。DP、TP、PP、EP、checkpoint、数据加载、推理请求各自产生不同流量形态，网络拓扑和调度系统必须理解这些差异。

\subsection{Oversubscription}
\label{subsec:oversubscription}

Oversubscription 指下行接入带宽大于上行汇聚带宽。例如一个 leaf switch 下接入 16 台服务器，每台 400Gb/s，总下行 6.4Tb/s，但上行到 spine 总带宽只有 3.2Tb/s，则 oversubscription ratio 为 2:1。普通数据中心为了降低成本，经常接受一定 oversubscription。但 AI 训练网络对 oversubscription 非常敏感。

如果多个节点之间频繁进行 all-reduce 或 all-to-all，跨 leaf 流量会占用上行带宽。oversubscription 会导致：

\begin{itemize}
  \item collective 带宽下降；
  \item 某些训练 step 长尾变长；
  \item 多作业混部互相影响；
  \item RoCE PFC/ECN 更容易触发；
  \item NCCL timeout 风险增加。
\end{itemize}

因此，大规模训练集群通常倾向于低 oversubscription 或 non-blocking fabric，尤其是训练网络。推理集群则可以根据 workload 做分层设计：prefill-heavy、decode-heavy、embedding-heavy、存储访问等流量可能使用不同网络或 QoS 策略。

\subsection{East-West Traffic}
\label{subsec:east-west-traffic}

AI 集群中的主要流量通常是 east-west traffic，即节点之间的横向流量，而不是 north-south traffic。传统 Web 服务中，请求可能从用户进入负载均衡，再访问服务、缓存、数据库，南北向流量明显。而分布式训练中，大多数通信发生在训练节点之间：rank 与 rank 交换梯度、参数分片、token、activation。

East-west traffic 的特点是：

\begin{itemize}
  \item 流量同步性强，多个 rank 同时发起 collective；
  \item 数据量大，容易形成突发；
  \item 通信模式规则，但受并行策略影响；
  \item all-to-all 可能造成复杂多对多流量；
  \item 单个慢 rank 会拖慢整个 job。
\end{itemize}

这也是为什么 AI 集群调度应尽量把同一个训练作业放在网络拓扑上相邻、带宽充足的节点集合中，而不是随机散布在整个数据中心。后续第 20 章讨论 Kubernetes 调度时，我们会看到 topology-aware scheduling、gang scheduling 和 queueing system 如何与网络拓扑结合。

\subsection{网络监控指标}
\label{subsec:network-monitoring-metrics}

AI 集群网络必须可观测。只监控 GPU utilization 不够，只监控交换机端口 up/down 也不够。至少需要覆盖以下层次：

\begin{longtable}{p{3.2cm}p{5.0cm}p{5.2cm}}
  \caption{AI 集群网络常见监控指标}
  \label{tab:ai-network-monitoring} \\
  \toprule
  层次 & 指标 & 说明 \\
  \midrule
  \endfirsthead
  \toprule
  层次 & 指标 & 说明 \\
  \midrule
  \endhead
  NIC & TX/RX bandwidth & 网卡发送和接收带宽，判断 rail 是否均衡 \\
  NIC & packet error / retry & 发现链路错误、丢包、重传和 RoCE 异常 \\
  NIC & RDMA counters & 包括 RDMA read/write/send、completion error 等 \\
  PCIe & GPU-NIC 拓扑和吞吐 & 判断是否跨 NUMA、跨 switch 或 P2P 不可用 \\
  Switch & port utilization & 发现热点链路和 oversubscription \\
  Switch & buffer occupancy & 判断是否存在队列积压和拥塞 \\
  Switch & PFC pause frame & 发现 PFC 触发频率、pause storm 或拥塞扩散 \\
  Switch & ECN marking & 判断拥塞标记是否合理，是否过早或过晚 \\
  NCCL & algbw / busbw & 验证 collective 有效带宽 \\
  NCCL & timeout / retry / fallback & 发现网络异常、fallback 到 socket 或连接问题 \\
  Job & step time / comm time & 判断通信是否在训练关键路径 \\
  Job & straggler rank & 找出拖慢整个作业的慢 rank \\
  \bottomrule
\end{longtable}

网络监控的关键是建立 baseline。同一类节点、同一类网卡、同一类拓扑，应有稳定的 NCCL tests 结果和端口带宽曲线。一旦某台机器或某条链路偏离 baseline，就应自动标记为可疑节点，避免调度大作业。

\section{实践：RDMA / RoCE 排查流程}
\label{sec:rdma-roce-debug-workflow}

本节给出一个常见排查流程。假设一个多节点训练作业扩展性异常，单机 8 卡正常，两机 16 卡明显变慢，或者 NCCL 偶发 timeout。

\subsection{第一步：确认 NCCL 是否走预期网络}
\label{subsec:debug-nccl-network}

首先打开 NCCL 日志：

\begin{lstlisting}[language=bash,caption={打开 NCCL 网络调试日志},label={lst:nccl-net-debug}]
export NCCL_DEBUG=INFO
export NCCL_DEBUG_SUBSYS=INIT,NET,GRAPH,TOPO,COLL
\end{lstlisting}

观察是否出现以下问题：

\begin{itemize}
  \item 没有识别到 IB/RoCE 设备；
  \item fallback 到 Socket；
  \item 使用了错误网卡；
  \item GID index 不正确；
  \item GPUDirect RDMA 未启用；
  \item 多 rail 未被利用；
  \item 某些 rank 连接失败或 timeout。
\end{itemize}

\subsection{第二步：运行基础带宽测试}
\label{subsec:debug-bandwidth-test}

在训练框架之外运行 NCCL tests：

\begin{lstlisting}[language=bash,caption={多节点 NCCL all-reduce 基线测试},label={lst:multi-node-nccl-baseline}]
mpirun -np 16 -H node0:8,node1:8 \
  -x NCCL_DEBUG=INFO \
  -x NCCL_IB_HCA=mlx5_0,mlx5_1 \
  ./build/all_reduce_perf -b 8M -e 8G -f 2 -g 1
\end{lstlisting}

如果 NCCL tests 已经异常，说明问题在通信基础设施或配置，不要先调模型。应继续检查 HCA、链路、交换机和系统配置。

\subsection{第三步：检查 HCA 与链路状态}
\label{subsec:debug-hca-link}

常见命令包括：

\begin{lstlisting}[language=bash,caption={检查 RDMA 设备和链路状态的常用命令},label={lst:rdma-link-check}]
# 查看 RDMA 设备
ibv_devinfo

# 查看网卡和端口
ibstat

# 查看 RDMA link
rdma link show

# 查看网卡统计
ethtool -S <interface>

# 查看 IP 与接口
ip addr show
\end{lstlisting}

检查内容包括：

\begin{itemize}
  \item link 是否 up；
  \item 速率是否符合预期；
  \item MTU 是否一致；
  \item 是否存在大量 error、drop、retry；
  \item RoCE 接口 IP、GID、VLAN、priority 是否正确；
  \item 多 NIC 是否全部可用。
\end{itemize}

\subsection{第四步：检查 RoCE QoS、PFC 与 ECN}
\label{subsec:debug-roce-qos}

RoCE 环境中，需要确认主机和交换机配置一致：

\begin{itemize}
  \item RoCE 流量是否映射到正确 priority；
  \item PFC 是否只对目标 priority 启用；
  \item ECN 阈值是否合理；
  \item DSCP/PCP 是否在交换机路径中被保留；
  \item 是否存在 pause frame 暴增；
  \item 是否存在 ECN marking 异常；
  \item 是否有普通业务流量混入 RoCE 无损队列。
\end{itemize}

如果 PFC pause frame 持续增加，同时训练 step 出现长尾，说明可能存在拥塞扩散或队列阻塞。此时需要结合交换机 telemetry 和作业调度信息，判断是否有多个大作业共享同一瓶颈链路。

\subsection{第五步：检查 GPU-NIC 拓扑}
\label{subsec:debug-gpu-nic-topology}

运行：

\begin{lstlisting}[language=bash,caption={检查 GPU 与 NIC 拓扑},label={lst:gpu-nic-topology}]
nvidia-smi topo -m
\end{lstlisting}

重点观察：

\begin{itemize}
  \item 每张 GPU 到每张 NIC 的路径；
  \item 是否跨 NUMA；
  \item 是否跨 PCIe switch；
  \item 多 rail 是否与 GPU 分组匹配；
  \item local rank 是否绑定到了合适 GPU。
\end{itemize}

如果某些 rank 使用远端 NIC，跨节点带宽可能明显下降。可以通过调整 \texttt{CUDA\_VISIBLE\_DEVICES}、launcher rank mapping、NCCL HCA 选择或调度器节点分配策略改善。

\subsection{第六步：结合训练 Timeline}
\label{subsec:debug-training-timeline}

最后，用 Nsight Systems 或框架 profiler 观察真实训练 step。需要判断：

\begin{itemize}
  \item NCCL kernel 是否在关键路径；
  \item backward 与 all-reduce 是否重叠；
  \item 是否有某个 rank 明显慢；
  \item 通信前后是否存在 CPU 空洞；
  \item checkpoint 或 dataloader 是否与训练通信争用网络；
  \item MoE all-to-all 是否造成突发拥塞。
\end{itemize}

只有把网络指标与训练 timeline 对齐，才能避免误判。例如 NCCL tests 正常但训练慢，可能是框架 bucket 太小、overlap 不好、rank mapping 不合理或其他 I/O 干扰；训练 timeline 中 NCCL 长尾明显，同时交换机 PFC/ECN 异常，则更可能是网络拥塞。

\section{设计原则与经验法则}
\label{sec:rdma-roce-design-principles}

\subsection{原则一：训练网络与普通业务网络分离}
\label{subsec:principle-separate-network}

大规模训练流量具有突发、大带宽、同步等待和长尾敏感特征。把训练 RDMA/RoCE 流量与普通业务流量混在同一无隔离网络中，容易互相影响。实践中应尽量使用独立训练网络，或至少通过 QoS、VLAN、VRF、优先级队列和带宽隔离明确区分。

\subsection{原则二：先做无损，再做拥塞控制}
\label{subsec:principle-lossless-congestion}

RoCE 的基础是避免丢包，但只靠 PFC 并不够。PFC 可以阻止队列溢出，却可能造成拥塞扩散。ECN 和端到端拥塞控制用于在不丢包的情况下调节发送速率。一个稳定 RoCE 网络需要 PFC、ECN、QoS、buffer 和路由共同配置。

\subsection{原则三：拓扑感知比盲目加带宽更重要}
\label{subsec:principle-topology-aware-network}

多 NIC、多 rail、多 leaf 的网络中，平均带宽不代表实际作业带宽。一个错误的 rank mapping 或 NIC 绑定可能让通信集中在少数链路。训练调度器应该理解节点、GPU、NIC、leaf、spine 的关系，把同一作业放在拓扑上更紧凑、更均衡的位置。

\subsection{原则四：用 NCCL Tests 建立集群基线}
\label{subsec:principle-nccl-baseline}

每个 AI 集群上线前都应建立 NCCL tests 基线，包括单机、双机、跨 leaf、多机规模下的 all-reduce、all-gather、reduce-scatter、all-to-all。后续每次驱动升级、交换机配置变更、NCCL 版本升级、系统镜像更新后，都应重新验证基线。

\subsection{原则五：监控要覆盖端到端路径}
\label{subsec:principle-end-to-end-monitoring}

只监控 GPU 不够，只监控交换机也不够。一次跨节点 GPU 通信跨越 GPU HBM、PCIe、NIC、交换机、对端 NIC、对端 PCIe、对端 GPU HBM。任何一段异常都可能拖慢训练。监控和告警应覆盖端到端路径，并能关联到具体 job、rank、node、NIC 和 switch port。

\section{本章小结}
\label{sec:rdma-roce-summary}

本章从多机 AI 集群的角度介绍了 RDMA、RoCE、InfiniBand、GPUDirect RDMA 与集群网络设计。

\begin{itemize}
  \item 多机训练依赖高性能网络，因为数据并行、张量并行、流水线并行、ZeRO 和 MoE 都会产生跨节点通信。
  \item 跨节点 all-reduce、reduce-scatter、all-gather 和 all-to-all 的性能直接影响训练 step time 和扩展效率。
  \item RDMA 通过 kernel bypass、zero-copy、queue pair、completion queue 和 memory registration 降低延迟与 CPU 开销。
  \item InfiniBand 是面向高性能 RDMA 的专用网络，RoCE 则把 RDMA 承载在以太网上，便于复用以太网生态。
  \item RoCE 网络需要重点关注 lossless Ethernet、PFC、ECN、QoS、拥塞控制和交换机 buffer，否则少量拥塞就可能导致训练长尾或 NCCL timeout。
  \item GPUDirect RDMA 允许 NIC 直接访问 GPU memory，减少 CPU host memory staging，是跨节点 GPU collective 的关键能力。
  \item 多 NIC、多 rail 节点中，GPU-NIC 亲和性、rank mapping 和 NCCL HCA 选择会显著影响实际带宽。
  \item AI 集群网络设计需要考虑 fat-tree、dragonfly、oversubscription、east-west traffic 和网络可观测性。
  \item 排查网络问题应从 NCCL 日志、NCCL tests、HCA 状态、RoCE QoS、GPU-NIC 拓扑和训练 timeline 多层结合入手。
\end{itemize}

至此，我们已经从单个 GPU、单机多卡互联，扩展到了多机 RDMA/RoCE 网络。下一章将进入集群资源调度层：当一个 AI 集群同时运行训练、推理、评测、数据处理和多租户任务时，Kubernetes 如何管理 GPU？Device Plugin、MIG、gang scheduling、Volcano、Kueue、拓扑感知调度和 AI 集群可观测性又如何协同工作？这些内容将进入第 20 章：Kubernetes 在 AI 集群中的调度。
